%=============================================================================
\chapter{Vanilla BART}
\label{ch:vanilla}

This chapter derives the original BART of Chipman, George \& McCulloch (2010)
as a special case of the model in \autoref{ch:full}: univariate outcome,
homoskedastic Gaussian noise.

%-----------------------------------------------------------------------------
\section{Substitutions}
\label{sec:subst}

The specialisation is
\begin{equation}
	\label{eq:subst}
	\begin{alignedat}{4}
		\setB &= \emptyset,  & \qquad k     &= 1,             & \qquad w     &\equiv 1,      & \qquad \vs  &= (1/p, \dots, 1/p), \\
		\va_i &\equiv 1,     & \qquad \mLam &= 1/\sigsq,      & \qquad \mLamm &= 1/\smsq,    & \qquad \Psi &= \nu\lambda, \\
		\nmindn &= 0,        & \qquad \nminleaf &= 5 .
	\end{alignedat}
\end{equation}
All $k \times k$ matrices become scalars, there are no latent outcomes
($\vyt_i = \vy_i$), the split-variable probabilities $\vs$ are fixed rather
than drawn from \eqref{eq:mv-s-prior} (so $\theta$ and its hyperprior play
no role), and the dictionary between the two chapters is:
\begin{center}
	\begin{tabular}{lll}
		\toprule
		quantity                & Chapter~\ref{ch:full}                                           & this chapter                                           \\
		\midrule
		outcome                 & $\vy_i \in \R^k$                                                & $y_i \in \R$                                           \\
		leaf parameter          & $\vmu_\ell \in \R^k$                                            & $\mu_\ell \in \R$                                      \\
		leaf prior              & $\N_k(\mathbf{0}, \mLamm^{-1})$                                 & $\N(0, \smsq)$                                         \\
		error precision\ at $i$ & $\mA_i\,\mLam\,\mA_i = \mLam\odot\tilde\mA_i$                   & $1/\sigsq$                                             \\
		error prec.\ prior      & $\W(\nu, \Psi)$                                                 & $\Gamma(\nu/2,\nu\lambda/2)$                           \\
		per-leaf count          & $\mN_\ell = \sum_{i\in I_\ell}\tilde\mA_i$                      & $n_\ell$                                               \\
		per-leaf score          & $S_\ell = \sum_{i\in I_\ell}\tilde\mA_i\odot(\Ind_k\vr_i^\top)$ & $S_\ell = \sum_{i\in I_\ell} r_i$                      \\
		precision sum           & $\mP_\ell = \mLam\odot\mN_\ell$                                 & $n_\ell/\sigsq$                                        \\
		score vector            & $\vB_\ell = (\mLam\odot S_\ell)\,\Ind_k$                        & $S_\ell/\sigsq$                                        \\
		leaf posterior          & $\N_k\!\bigl(V_\ell\vB_\ell, V_\ell\bigr)$                      & $\N(S_\ell V_\ell/\sigsq, V_\ell)$                     \\
		~ post.\ precision      & $\mLamm + \mP_\ell$                                             & $n_\ell/\sigsq + 1/\smsq$                              \\
		log-det term            & $\frac{1}{2}\log|\mLamm| - \frac{1}{2}\log|\mLamm+\mP_\ell|$  & $\frac{1}{2}\log\!\frac{\sigsq}{n_\ell\smsq+\sigsq}$ \\
		quadratic term          & $\frac{1}{2}\vB_\ell^\top(\mLamm+\mP_\ell)^{-1}\vB_\ell$       & $\frac{\smsq S_\ell^2}{2\sigsq(n_\ell\smsq+\sigsq)}$  \\
		$\mLam$ post.\ rate     & $\Psi + \sum_i (\va_i\odot\veps_i)(\va_i\odot\veps_i)^\top$     & $\nu\lambda + \sum_i \veps_i^2$                        \\
		\bottomrule
	\end{tabular}
\end{center}
All tree-prior and proposal-kernel factors are unchanged, but for the
uniform variable choice $\vs \equiv 1/p$, the uniform leaf choice
$w \equiv 1$, and the count thresholds: no decision-node threshold, and the
CGM requirement of at least five training points per leaf.

%-----------------------------------------------------------------------------
\section{Notation and data}
\label{sec:notation}

We observe $n$ pairs $(\vx_i, y_i)_{i=1}^n$ with $\vx_i \in \R^p$ and
$y_i \in \R$. Stack the responses into $\vy = (y_1, \dots, y_n)^\top$ and the
covariates into $X \in \R^{n \times p}$ as in \S\ref{sec:mv-notation}. The
precision scales are all $\va_i = 1$ and are dropped from the notation.

Trees $T$, leaves $\ell \in L(T)$, and the assignment $\ell(\vx; T)$ are
defined exactly as in \S\ref{sec:mv-notation}. The leaf array is now a set
of scalars,
\[
	\vM = \{ \mu_\ell : \ell \in L(T) \} \subset \R,
\]
with regression function and sum of trees
\begin{equation}
	\label{eq:sum-of-trees}
	g(\vx; T, \vM) = \sum_{\ell \in L(T)} \mu_\ell\, \Ind[\ell(\vx; T) = \ell],
	\qquad
	f(\vx) = \sum_{t=1}^m g(\vx;\, T_t, \vM_t).
\end{equation}

%-----------------------------------------------------------------------------
\section{The generative model}
\label{sec:model}

\subsection{Conditional outcome distribution}

With $k = 1$ and $\mA_i = 1$, the precision $\mA_i\mLam\mA_i$ in
\eqref{eq:mv-lik} is $1/\sigsq$ for every $i$, so
\begin{equation}
	\label{eq:lik}
	y_i \mid T_{1:m}, \vM_{1:m}, \sigsq \stackrel{\text{ind}}{\sim}
	\N\!\left(f(\vx_i), \sigsq\right),
	\qquad i = 1, \dots, n.
\end{equation}

\subsection{Tree-structure prior}

As in \S\ref{sec:mv-tree-prior}, with $\vs$ fixed at $(1/p, \dots, 1/p)$: a
node at depth $d$ is internal with probability
$p_{\mathrm{NT}}(d) = \alpha(1+d)^{-\beta}$ of \eqref{eq:pnt}, truncated at
$d_{\max}$, and given that it is internal its splitting rule $(v, c)$ is
drawn in two stages, the variable $v$ uniformly over the eligible set
$E(\eta)$ by \eqref{eq:mv-var-prior}, then the cutpoint $c$ uniformly over
the cutpoints of $v$ compatible with the ancestor splits. Since the number
of compatible cutpoints varies with $v$, this is not uniform over the
admissible $(v, c)$ pairs. The support is restricted to trees with at least
$\nminleaf = 5$ training points in every leaf, the constraint of CGM (2010);
there is no decision-node threshold, $\nmindn = 0$. bartz's \texttt{gbart}
sets $\nminleaf = 5$ but keeps $\nmindn = 10$: with
$\nmindn = 2\nminleaf$ the leaf threshold already implies the decision-node
one, so the support, hence the target, is the same.

\subsection{Leaf prior}

Equation \eqref{eq:mv-leaf-prior} with $\mLamm = 1/\smsq$ is
\begin{equation}
	\label{eq:leaf-prior}
	\mu_\ell \mid T_t \stackrel{\text{i.i.d.}}{\sim}
	\N\!\left(0, \smsq\right),
	\qquad \ell \in L(T_t).
\end{equation}
The variance $\smsq$ is calibrated so that the prior on the sum of trees
$f(\vx)$ has standard deviation $\sqrt{m}\,\sm$ controlled by the user. Bartz
exposes the inverse, $1/\smsq$, via the field
\texttt{Forest.leaf\_prior\_cov\_inv}; stochtree's
\texttt{GaussianConstantLeafModel} stores $\tau \equiv \smsq$.

\subsection{Error variance prior}
\label{sec:sigma-prior}

For $k = 1$ the Wishart prior \eqref{eq:mv-iw-prior} on the scalar
precision $\mLam = 1/\sigsq$ has density proportional to
$\mLam^{(\nu - 2)/2}\exp(-\Psi\mLam/2)$, which is the gamma distribution
$\Gamma(\nu/2,\, \Psi/2)$ in the shape--rate parametrisation. Equivalently,
the variance follows the conjugate inverse-gamma prior
\begin{equation}
	\label{eq:sigma-prior}
	\sigsq \sim \IG\!\left(\frac{\nu}{2}, \frac{\nu\lambda}{2}\right),
	\qquad \Psi = \nu\lambda,
\end{equation}
which in \texttt{bartz} is the \texttt{Wishart} object passed as
\texttt{error\_cov\_inv} to \texttt{mcmcstep.init}, with
\texttt{nu}~$=\nu$ and \texttt{rate}~$=\nu\lambda$. The user-facing
\texttt{sigma\_df} and \texttt{sigma\_scale} arguments of \texttt{bartz.Bart}
set $\nu$ and $\lambda = \texttt{sigma\_scale}^2$.

\subsection{Bayesian network}
\label{sec:bn}

With $\vs$ fixed, the factors $p(\theta)\,p(\vs \mid \theta)$ of
\eqref{eq:mv-joint} disappear and $p(T_t \mid \vs)$ becomes a fixed tree
prior $p(T_t)$, so the joint density reads
\begin{equation}
	\label{eq:joint}
	p(\vy, T_{1:m}, \vM_{1:m}, \sigsq)
	= p(\sigsq) \prod_{t=1}^m p(T_t)\, p(\vM_t \mid T_t)
	\cdot \prod_{i=1}^n p(y_i \mid T_{1:m}, \vM_{1:m}, \sigsq),
\end{equation}
with $p(\vM_t \mid T_t) = \prod_{\ell \in L(T_t)} \N(\mu_\ell;\, 0, \smsq)$.
The graphical model is that of \S\ref{sec:mv-bn} without the $\theta$ and
$\vs$ nodes and without the $\va_i$ input node:
\begin{center}
	\begin{tikzpicture}[node distance=0.9cm and 1.6cm,
			every node/.style={font=\footnotesize}]
		% Top row: hyperparameters
		\node[const]                       (alpha) {$\alpha,\beta$};
		\node[const, right=2.2cm of alpha] (sm)    {$\smsq$};
		\node[const, right=2.2cm of sm]    (nu)    {$\nu,\lambda$};
		% Middle row: latent variables
		\node[latent, below=of alpha] (T)  {$T_t$};
		\node[latent, below=of sm]    (mu) {$\mu_\ell^{(t)}$};
		\node[latent, below=of nu]    (s2) {$\sigsq$};
		% Bottom row: data
		\node[obs, below=1.6cm of mu] (y) {$y_i$};
		\node[const, left=of y]       (x) {$\vx_i$};

		\edge {alpha} {T};
		\edge {sm}    {mu};
		\edge {nu}    {s2};
		\edge {T}     {mu};
		\edge {mu}    {y};
		\edge {s2}    {y};
		\edge {x}     {y};
		\draw[->, >={triangle 45}] (T) to[bend right=25] (y);

		\plate {leaves} {(mu)}        {$\ell \in L(T_t)$};
		\plate {trees}  {(T)(leaves)} {$t = 1,\dots,m$};
		\plate {data}   {(x)(y)}      {$i = 1,\dots,n$};
	\end{tikzpicture}
\end{center}

%-----------------------------------------------------------------------------
\section{Single-tree decomposition and partial residuals}
\label{sec:partial-resid}

The backfitting decomposition of \S\ref{sec:mv-partial-resid} applies with
scalar partial residuals
\begin{equation}
	\label{eq:partial-resid}
	r_i^{(t)} \coloneqq y_i - \sum_{t' \ne t} g(\vx_i;\, T_{t'}, \vM_{t'}),
	\qquad
	r_i^{(t)} \mid T_t, \vM_t, \sigsq \stackrel{\text{ind}}{\sim}
	\N\!\left(g(\vx_i;\, T_t, \vM_t),\, \sigsq\right).
\end{equation}
As before we fix $t$ and write $r_i = r_i^{(t)}$, $\vr = \vr^{(t)}$,
$T = T_t$, $\vM = \vM_t$.

With $\va_i \equiv 1$ the outer product $\tilde\mA_i = \va_i\va_i^\top$ of
\eqref{eq:mv-Atilde} is the scalar $1$, so the matrix sufficient statistics
\eqref{eq:mv-suffstats-N}--\eqref{eq:mv-suffstats-S} collapse to a count and
a sum. For each leaf $\ell \in L(T)$,
\begin{equation}
	\label{eq:suffstats}
	I_\ell = \{\, i : \ell(\vx_i; T) = \ell \,\},
	\qquad
	\mN_\ell = n_\ell = |I_\ell|,
	\qquad
	S_\ell = \sum_{i \in I_\ell} r_i,
	\qquad
	Q_\ell = \sum_{i \in I_\ell} r_i^{2},
\end{equation}
where $Q_\ell$ is introduced only to write the full marginal density
explicitly. The $\mLam$-dependent quantities
\eqref{eq:mv-Pi}--\eqref{eq:mv-suffstats-B} become
\begin{equation}
	\label{eq:suffstats-lam}
	\mP_i = \frac{1}{\sigsq},
	\qquad
	\mP_\ell = \frac{n_\ell}{\sigsq},
	\qquad
	\vB_\ell = \frac{S_\ell}{\sigsq}.
\end{equation}

The bartz code accumulates $S_\ell$ and $n_\ell$ in
\texttt{sum\_resid} / \texttt{compute\_count\_trees}; stochtree calls them
\texttt{sum\_yw} (with weight $1$) and \texttt{n} in
\texttt{GaussianConstantSuffStat}.

%-----------------------------------------------------------------------------
\section{Full conditional of the leaves}
\label{sec:full-cond}

\begin{prop}[Leaf full conditional, univariate homoskedastic]
	\label{prop:leaf}
	For each $\ell \in L(T)$,
	\begin{equation}
		\label{eq:leaf-fc}
		\mu_\ell \mid T, \sigsq, \vr \sim
		\N\!\left(
		\underbrace{\frac{S_\ell/\sigsq}{n_\ell/\sigsq + 1/\smsq}}_{=\,\hat\mu_\ell}
		,\,
		\underbrace{\frac{1}{n_\ell/\sigsq + 1/\smsq}}_{=\,V_\ell}
		\right),
	\end{equation}
	and the leaves are conditionally independent across $\ell$.
\end{prop}

\begin{proof}
	Substitute \eqref{eq:suffstats-lam} and $\mLamm = 1/\smsq$ into
	Proposition~\ref{prop:mv-leaf}: the posterior precision is
	$\mLamm + \mP_\ell = 1/\smsq + n_\ell/\sigsq$, so
	$V_\ell = (\mLamm + \mP_\ell)^{-1}$ and
	$\hat\mu_\ell = V_\ell\vB_\ell = V_\ell\,S_\ell/\sigsq$ are as stated.
\end{proof}

\paragraph{Equivalent form.}
Multiplying numerator and denominator of $\hat\mu_\ell$ by $\sigsq\smsq$ we
recover the closed form used in stochtree's
\texttt{PosteriorParameterMean} / \texttt{PosteriorParameterVariance}:
\begin{equation}
	\label{eq:leaf-fc-tau}
	\hat\mu_\ell = \frac{\smsq\, S_\ell}{n_\ell\,\smsq + \sigsq},
	\qquad
	V_\ell = \frac{\smsq\,\sigsq}{n_\ell\,\smsq + \sigsq}.
\end{equation}
The bartz code uses the precision form: in
\texttt{\_precompute\_leaf\_terms\_uv} it computes
\texttt{var\_post}~$= V_\ell$ and \texttt{mean\_factor}~$= V_\ell / \sigsq$,
then samples $\mu_\ell = \texttt{mean\_factor} \cdot S_\ell + \sqrt{V_\ell}\, Z$
with $Z \sim \N(0,1)$ (the draw $\sqrt{V_\ell} Z$ is precomputed as
\texttt{centered\_leaves}, the product with $S_\ell$ happens in
\texttt{accept\_move\_and\_sample\_leaves}).

%-----------------------------------------------------------------------------
\section{Marginal density with leaves integrated out}
\label{sec:marg-lik}

Write $\vr_\ell = (r_i)_{i \in I_\ell}$ for the vector of partial residuals
falling in leaf $\ell$ (the $k = 1$ case of $\mR_\ell$ in
\eqref{eq:mv-Rell}). Conditional on $T$ and $\sigsq$,
\[
	\vr_\ell \mid T, \sigsq \sim
	\N_{n_\ell}\!\left(\mathbf{0},\, \sigsq I_{n_\ell} + \smsq\,\Ind\Ind^\top\right),
\]
the leaves being independent across $\ell$.

\begin{prop}[Per-leaf marginal density, univariate homoskedastic]
	\label{prop:marglik}
	For each leaf $\ell$,
	\begin{equation}
		\label{eq:marglik-full}
		\log p(\vr_\ell \mid T, \sigsq)
		= -\frac{n_\ell}{2}\log(2\pi)
		- \frac{n_\ell}{2}\log\sigsq
		+ \frac{1}{2}\log\frac{\sigsq}{n_\ell\,\smsq + \sigsq}
		- \frac{Q_\ell}{2\sigsq}
		+ \frac{\smsq\, S_\ell^{2}}{2\sigsq\,(n_\ell\,\smsq + \sigsq)},
	\end{equation}
	and the joint log marginal of the tree is
	\(
	\log p(\vr \mid T, \sigsq) = \sum_{\ell \in L(T)} \log p(\vr_\ell \mid T, \sigsq).
	\)
\end{prop}

\begin{proof}
	Specialise the three lines of Proposition~\ref{prop:mv-marglik} in turn.
	Line \eqref{eq:mv-marglik-a}: with $k = 1$ and
	$|\mA_i\mLam\mA_i| = 1/\sigsq$,
	\[
		-\frac{n_\ell}{2}\log(2\pi)
		+ \frac{1}{2}\sum_{i \in I_\ell}\log\frac{1}{\sigsq}
		= -\frac{n_\ell}{2}\log(2\pi) - \frac{n_\ell}{2}\log\sigsq .
	\]
	Line \eqref{eq:mv-marglik-b}: with $\mLamm = 1/\smsq$ and
	$\mLamm + \mP_\ell = 1/\smsq + n_\ell/\sigsq
		= (n_\ell\,\smsq + \sigsq)/(\sigsq\smsq)$,
	\[
		\frac{1}{2}\log\frac{1}{\smsq}
		- \frac{1}{2}\log\frac{n_\ell\,\smsq + \sigsq}{\sigsq\smsq}
		= \frac{1}{2}\log\frac{\sigsq}{n_\ell\,\smsq + \sigsq} .
	\]
	Line \eqref{eq:mv-marglik-c}: with $\mP_i = 1/\sigsq$ and
	$\vB_\ell = S_\ell/\sigsq$,
	\[
		-\frac{1}{2}\sum_{i \in I_\ell}\frac{r_i^2}{\sigsq}
		+ \frac{1}{2}\,\frac{S_\ell^2}{\sigma^4}\,
		\frac{\sigsq\smsq}{n_\ell\,\smsq + \sigsq}
		= -\frac{Q_\ell}{2\sigsq}
		+ \frac{\smsq\, S_\ell^{2}}{2\sigsq\,(n_\ell\,\smsq + \sigsq)} .
	\]
	Adding the three lines gives \eqref{eq:marglik-full}.
\end{proof}

\paragraph{Reduced form for tree-structure moves.}
The terms $-\frac{n_\ell}{2}\log(2\pi)$, $-\frac{n_\ell}{2}\log\sigsq$ and
$-Q_\ell/(2\sigsq)$ are the specialisation of the additive terms singled out
in \S\ref{sec:mv-marg-lik}: they cancel when a node $\eta$ is compared
against its split $(\eta_L, \eta_R)$, and only the data-dependent
log-determinant and quadratic-form terms remain:
\begin{equation}
	\label{eq:reduced-marglik}
	\log p(\vr_\ell \mid T, \sigsq) =
	\underbrace{\frac{1}{2}\log\frac{\sigsq}{n_\ell\,\smsq + \sigsq}}_{\text{log-det term}}
	+ \underbrace{\frac{\smsq\, S_\ell^{2}}{2\sigsq\,(n_\ell\,\smsq + \sigsq)}}_{\text{quadratic term}}
	+ (\text{cancels in splits / prunes}).
\end{equation}
This is exactly the quantity returned by
\texttt{GaussianConstantLeafModel::SplitLogMarginalLikelihood} and
\texttt{NoSplitLogMarginalLikelihood} in stochtree, with $\tau = \smsq$ and
\texttt{global\_variance} $= \sigsq$.

%-----------------------------------------------------------------------------
\section{Acceptance ratios for GROW and PRUNE}
\label{sec:lkratios}

\subsection{The GROW move}
\label{sec:grow}

The GROW proposal and its acceptance ratio
\begin{equation}
	\label{eq:mh}
	\mathrm{MH}_{\mathrm{grow}}
	= \frac{q(T \mid T')}{q(T' \mid T)}
	\cdot \frac{p(T')}{p(T)}
	\cdot \frac{p(\vr \mid T', \sigsq)}{p(\vr \mid T, \sigsq)}
\end{equation}
are those of \S\ref{sec:mv-grow}. The non-likelihood factor (transition
ratio times prior ratio) does not touch the leaf or noise model, so it is
\eqref{eq:mv-nonlk-grow} with the uniform leaf choice $w \equiv 1$:
\begin{equation}
	\label{eq:nonlk-grow}
	\frac{q(T \mid T')\,p(T')}{q(T' \mid T)\,p(T)}
	= \frac{P_{\mathrm{prune}}(T')}{P_{\mathrm{grow}}(T)} \cdot
	\frac{|G(T)|}{|P(T')|} \cdot
	\frac{p_{d_\eta}}{1 - p_{d_\eta}} \cdot
	(1 - p_{d_\eta+1})^{2},
\end{equation}
with the same convention for non-growable children. Since $\nmindn = 0$,
$G(T)$ is count-free: every leaf with an admissible split above the depth
limit is growable. The leaf threshold enters only through the veto of
\S\ref{sec:mv-grow}: a GROW whose children would have fewer than five points
is rejected.

\paragraph{Marginal density ratio.}
Substituting \eqref{eq:suffstats-lam} into \eqref{eq:mv-lkratio-grow}, the
four log-determinants combine as in the proof of
Proposition~\ref{prop:marglik} and the three quadratic forms become
$\smsq S_*^2/[\sigsq(n_*\smsq + \sigsq)]$, giving
\begin{align}
	\label{eq:lkratio-grow}
	\log\frac{p(\vr \mid T', \sigsq)}{p(\vr \mid T, \sigsq)}
	 & = \frac{1}{2}\log\frac{\sigsq\,(n_\eta\,\smsq + \sigsq)}
	                    {(n_{\eta_L}\,\smsq + \sigsq)\,(n_{\eta_R}\,\smsq + \sigsq)}
	\notag                                                                              \\
	 & \quad + \frac{\smsq}{2\sigsq}\left[
		                                \frac{S_{\eta_L}^{2}}{n_{\eta_L}\,\smsq + \sigsq}
		                                + \frac{S_{\eta_R}^{2}}{n_{\eta_R}\,\smsq + \sigsq}
		                                - \frac{S_\eta^{2}}{n_\eta\,\smsq + \sigsq}
		                                \right].
\end{align}
Recall $n_\eta = n_{\eta_L} + n_{\eta_R}$ and $S_\eta = S_{\eta_L} + S_{\eta_R}$.

This is exactly what \texttt{bartz} computes in
\texttt{\_compute\_likelihood\_ratio\_uv}, from the quantities precomputed by
\texttt{\_precompute\_likelihood\_terms\_uv} for the three nodes
$* \in \{\eta_L, \eta_R, \eta\}$:
\begin{align*}
	\texttt{prelkv.lrt}_*       & = \frac{\texttt{mean\_factor}_*}{\sigsq}
	= \frac{\sigma^{-4}}{n_*/\sigsq + 1/\smsq}
	= \frac{\smsq}{\sigsq (n_* \smsq + \sigsq)},                                    \\
	\texttt{prelkv.log\_sqrt\_term} & = \frac{1}{2} \log\frac{\texttt{lrt}_{\eta_L} \cdot \texttt{lrt}_{\eta_R}}{\texttt{lrt}_0 \cdot \texttt{lrt}_\eta},
	\qquad \texttt{lrt}_0 = \frac{\smsq}{\sigma^4} \text{ (the value at $n = 0$)},
\end{align*}
so that \eqref{eq:lkratio-grow} equals
$\texttt{log\_sqrt\_term} + \frac{1}{2}\left(
	S_{\eta_L}^2 \cdot \texttt{lrt}_{\eta_L}
	+ S_{\eta_R}^2 \cdot \texttt{lrt}_{\eta_R}
	- S_\eta^2 \cdot \texttt{lrt}_\eta\right)$.
It is also algebraically identical to
\texttt{SplitLogMarginalLikelihood} $-$ \texttt{NoSplitLogMarginalLikelihood}
in stochtree (after substituting $\tau = \smsq$ and \texttt{sum\_w}~$= n_\ell$,
\texttt{sum\_yw}~$= S_\ell$).

\subsection{The PRUNE move}
\label{sec:prune}

As in \S\ref{sec:mv-prune}, the PRUNE acceptance ratio is the reciprocal of
the GROW ratio for the inverse move, so its log marginal-density part is the
negative of \eqref{eq:lkratio-grow} evaluated with
$(n_{\eta_L}, S_{\eta_L}, n_{\eta_R}, S_{\eta_R})$ read off the
to-be-deleted children, and its non-likelihood part is the reciprocal of
\eqref{eq:nonlk-grow} for the inverse move.

\subsection{Sampling the leaves after acceptance}

Once the structure update is finalised, the leaves of $T$ are resampled from
\eqref{eq:leaf-fc} and the running residuals refreshed as in
\S\ref{sec:mv-sample-leaves}.

%-----------------------------------------------------------------------------
\section{Updating the error variance}
\label{sec:sigma-update}

After all $m$ trees have been updated, Proposition~\ref{prop:mv-Sigma}(i) with
$\tilde{\bm{\veps}}_i = \veps_i = y_i - f(\vx_i)$ and $\Psi = \nu\lambda$
gives
\[
	\frac{1}{\sigsq} \mid \cdots \sim
	\W\!\left(\nu + n,\; \nu\lambda + \sum_{i=1}^n \veps_i^2\right)
	= \Gamma\!\left(\frac{\nu + n}{2},\; \frac{\nu\lambda + \sum_i \veps_i^2}{2}\right),
\]
by the same $k = 1$ identification as in \S\ref{sec:sigma-prior}. In terms
of the variance,
\begin{equation}
	\label{eq:sigma-fc}
	\sigsq \mid \cdots \sim
	\IG\!\left(
	\frac{\nu + n}{2},\;
	\frac{\nu\lambda + \sum_{i=1}^n (y_i - f(\vx_i))^2}{2}
	\right),
\end{equation}
which is the standard conjugate update for a Gaussian sample with known
mean and an inverse-gamma prior on the variance. Bartz draws it in
\texttt{\_step\_error\_cov\_inv\_diag} as a gamma variate for $1/\sigsq$
with shape $(\nu + n)/2$ and rate $(\nu\lambda + \sum_i \veps_i^2)/2$.

%-----------------------------------------------------------------------------
\section{Implementation pointers}
\label{sec:impl-vanilla}

\begin{description}
	\item[bartz]
	      \begin{itemize}
		      \item \texttt{src/bartz/mcmcstep/\_state.py}: state and priors.
		      \item \texttt{src/bartz/mcmcstep/\_moves.py}: GROW/PRUNE
		            proposals, \texttt{compute\_partial\_ratio}.
		      \item \texttt{src/bartz/mcmcstep/\_step.py}:
		            \texttt{complete\_ratio},
		            \texttt{\_precompute\_leaf\_terms\_uv},
		            \texttt{\_precompute\_likelihood\_terms\_uv},
		            \texttt{\_compute\_likelihood\_ratio\_uv},
		            \texttt{\_step\_error\_cov\_inv\_diag}. The univariate
		            case is dispatched to these by
		            \texttt{precompute\_leaf\_terms},
		            \texttt{precompute\_likelihood\_terms},
		            \texttt{compute\_likelihood\_ratio}, and
		            \texttt{step\_error\_cov\_inv}.
	      \end{itemize}
	\item[stochtree]
	      \begin{itemize}
		      \item \texttt{include/stochtree/leaf\_model.h}: closed-form
		            derivations in the doxygen header.
		      \item \texttt{src/leaf\_model.cpp}:
		            \texttt{GaussianConstantLeafModel::SplitLogMarginalLikelihood},
		            \texttt{NoSplitLogMarginalLikelihood},
		            \texttt{PosteriorParameterMean},
		            \texttt{PosteriorParameterVariance},
		            \texttt{SampleLeafParameters}.
		      \item \texttt{src/sampler.cpp}: acceptance / rejection logic.
	      \end{itemize}
\end{description}
